\chapter{Índice Maestro de Validaciones}
\label{ch:apendice_datos}

Este apéndice es el índice central de 28 verificaciones y métricas
reportadas. Cada fila relaciona el identificador con el origen, el script,
el resultado consignado y la evidencia principal en los Apéndices A o D.
Las referencias a figuras y tablas sirven para localizar el desarrollo; no
implican ejecuciones independientes ni cambian el alcance descrito junto a
cada evidencia.

{\small
\setlength{\tabcolsep}{3pt}
\renewcommand{\arraystretch}{1.03}
\sloppy
\begin{longtable}{c >{\raggedright\arraybackslash}p{2.1cm} >{\raggedright\arraybackslash}p{2.4cm} >{\raggedright\arraybackslash}p{2.0cm} >{\raggedright\arraybackslash}p{4.1cm} >{\raggedright\arraybackslash}p{2.2cm}}
\caption[Índice maestro de validaciones]{Índice de V01--V28 con scripts, valores recalculados o reportados y evidencia principal. El alcance indicado en cada fila distingue resultados reproducidos de históricos.}
\label{tab:28_validaciones} \\
\toprule
ID & Origen & Script asociado & Métrica & Valor consignado & Evidencia \\
\midrule
\endfirsthead
\caption[]{Índice maestro de validaciones (continuación)} \\
\toprule
ID & Origen & Script asociado & Métrica & Valor consignado & Evidencia \\
\midrule
\endhead
\midrule
\multicolumn{6}{r}{\textit{Continúa en la siguiente página...}} \\
\endfoot
\bottomrule
\endlastfoot

V01 & Strukov 2008 & \url{01_strukov_2008.py} & $R^2$ y ERP por rango & Estado/corriente $R^2=0.999657/0.996994$; ERP $=1.27/6.22\%$ & \url{validacion_strukov_metricas.csv}; Tabla~\ref{tab:metricas_validacion} \\
V02 & Strukov 2008 & \url{07_r_x_consistencia.py} & Error relativo $R(x)$ & $0.00\%$ reportado & Fig.~\ref{fig:07_r_x} \\
V03 & Strukov 2008 & \url{04_convergencia.py} & Orden de convergencia $p$ & $0.98$ reportado & Fig.~\ref{fig:09_dt_extremos} \\
V04 & Strukov 2008 & \url{09_dt_extremos.py} & Umbral máximo de $dt$ & $10^{-3}\,\mathrm{s}$ reportado & Fig.~\ref{fig:09_dt_extremos} \\
V05 & Strukov 2008 & \url{10_triangular.py} & Rango de $x(t)$ & $[0.0,0.38]$ reportado & Fig.~\ref{fig:val-triangular} \\
V06 & Prezioso 2014 & \url{02_prezioso_2014.py} & Histéresis I--V & Visual aprox.; sin métricas ni ramas identificadas & Fig.~\ref{fig:val-prezioso} \\
V07 & Modelo local inspirado en Prezioso 2014 & \url{22_prezioso_v2.py} & Salida I--V simulada; sin ajuste a datos experimentales & Sin ajuste cuantitativo & Fig.~\ref{fig:22_prezioso_v2} \\
V08 & Comparación cualitativa Prezioso 2015 S3a & \url{23_prezioso_s3.py} & Forma I--V sin puntos digitalizados & Cuali\-tativa & Fig.~\ref{fig:prezioso_s3a} \\
V09 & Modelo local & \url{23_prezioso_s3.py} & Matriz virtual $4\times4$ & Sintética; semilla 42 & Fig.~\ref{fig:prezioso_s3bc} \\
V10 & Modelo local & \url{23_prezioso_s3.py} & Histograma de matriz virtual & Derivado de V09 & Fig.~\ref{fig:prezioso_s3bc} \\
V11 & Modelo local HfO$_2$--LIF & \url{06_hfo2_isi.py} & ISI dentro de dos ciclos simulados & Resultado local reportado & Fig.~\ref{fig:hfo2_isi_hibrido} \\
V12 & Identidad interna del modelo & \url{08_rs_de_x.py} & $R_S$ frente a $x(t)$ & Verifi\-cación algebraica reportada & Fig.~\ref{fig:08_rs_x} \\
V13 & Wang 2025; referencia codificada & \url{13_wang_2025.py} & Parámetros de referencia & Resumen guardado; no es serie temporal & Tabla~\ref{tab:csv_datasets} \\
V14 & Jo 2010, atribuido & \url{24_jo2010_white.py} & STDP; holdout ($R^2$, RMSE) & $0.94220$; $2.13626$ pp; $n=4$; signo declarado & Tabla~\ref{tab:metricas_validacion} \\
V15 & Jo 2010 & \url{17_stdp_spikes.py} & Solapamiento de pulsos & $\Delta V_{\mathrm{pico}}=450\,\mathrm{mV}$ reportado & Fig.~\ref{fig:val-stdp-spikes} \\
V16 & Jo 2010 & \url{18_stdp_temporal.py} & Acumulación temporal & $R^2=0.99$ reportado & Fig.~\ref{fig:18_stdp_temporal} \\
V17 & Jo 2010 & \url{14_ltp_ltd.py} & Potenciación $\Delta G>0$ & Signo reportado & Fig.~\ref{fig:ltp_ltd_base} \\
V18 & Jo 2010 & \url{14_ltp_ltd.py} & Depresión $\Delta G<0$ & Signo reportado & Fig.~\ref{fig:ltp_ltd_base} \\
V19 & Jo 2010 & \url{16_ltp_ltd_ciclo.py} & Error de retorno (100 ciclos) & $0.84\%$ reportado & Fig.~\ref{fig:ltp_ltd_ciclo} \\
V20 & Jo 2010 & \url{20_vi_ltp_ltd.py} & Coherencia V-I-R-G & Error $=0.00\%$ reportado & Fig.~\ref{fig:20_vi_ltp_ltd} \\
V21 & Jo 2010 & \url{19_iv_histeresis.py} & Lazo pinzado en origen & $I(0)=0.0\,\mathrm{nA}$ reportado & CSV en Tabla~\ref{tab:csv_datasets} \\
V22 & Solución exacta & \url{03_lif_analitica.py} & $R^2$ subumbral LIF & 0.9999815 (archivado) & Fig.~\ref{fig:lif_fi_curve} \\
V23 & Variabilidad D2D & \url{05_d2d_c2c.py} & Test KS $p$-valor D2D & $0.8959$ reproducido & Fig.~\ref{fig:d2d_c2c_validation} \\
V24 & Variabilidad C2C sintética & \url{05_d2d_c2c.py} & Error relativo de varianza OU & $3.62\%$; prueba sintética & Fig.~\ref{fig:d2d_c2c_validation} \\
V25 & Ruido paramétrico & \url{11_ruido.py} & CV de $I_{pp}$ & $115\%$ reportado & Fig.~\ref{fig:11_ruido} \\
V26 & D2D/C2C & \url{12_estadistica.py} & KS $p$-valor & \shortstack{ON: 0.7092, OFF: 0.8351\\$x_0$: 0.6187, C2C: 0.0000} & Fig.~\ref{fig:12_estadistica} \\
V27 & Crossbar $1\times1$ & \url{21_crossbar_1x1.py} & Operación VMM ($I=G^{\mathsf T}V$) & Error $=0.00\%$ reportado & Fig.~\ref{fig:crossbar_1x1_completo} \\
V28 & Neurona LIF & \url{26_lif_fi_curve.py} & Curva $f-I$ frente a solución analítica & 14 corrientes; mismo modelo & Fig.~\ref{fig:lif_fi_curve} \\
\end{longtable}
}


% =====================================================================
\section{Registros del piloto comparativo REINFORCE/R-STDP}
\label{sec:anexo_piloto_rstdp}

Los resultados de aprendizaje por tarea se conservan fuera de la suite de
validaciones de dispositivos enumerada en la Tabla~\ref{tab:28_validaciones}.
El lote comparativo contiene 28 corridas: 24 condiciones para CartPole,
visi\'on y obst\'aculos (dos algoritmos, semillas 0 y 1, y $X_0=0.1$ y
$0.9$) y cuatro condiciones Car R-STDP. Cada carpeta contiene el resultado
JSON, la curva de {\'e}xito, la imagen del crossbar y matrices inicial y final.
El protocolo y el resumen est\'an en
\url{../../03_EVIDENCIAS/ENTRENAMIENTOS/COMPARATIVA_SEEDS_0_1_X0/}.

El criterio fue superar el 90\% de {\'e}xitos en una ventana m\'ovil de 300
episodios, equivalente a un m\'inimo de 271 {\'e}xitos; el presupuesto
m\'aximo fue 10,000 episodios. La definici\'on binaria de {\'e}xito var\'ia
por dominio y se describe en cada JSON. Los dos registros CartPole REINFORCE
incompletos no se clasifican como fallos al m\'aximo. Las corridas Car del
primer lote usaron evaluaci\'on cada 200 episodios; no deben combinarse como
comparaci\'on directa con la repetici\'on corregida, que us\'o cada 100.

La repetici\'on corregida incluye cuatro corridas Car R-STDP con semillas 2
y 4 y $X_0=0.1$ y $0.9$. La semilla 2 alcanz\'o el criterio en los episodios
750 y 1548; las dos condiciones de la semilla 4 llegaron a 10,000 episodios
sin alcanzarlo. Sus registros y resumen est\'an en
\url{../../03_EVIDENCIAS/ENTRENAMIENTOS/CAR_RSTDP_EVAL100_SEEDS_2_4_X0/}.
Estos pilotos documentan ejecuciones bajo configuraciones concretas; no
constituyen validaci\'on en hardware ni una estimaci\'on estad\'istica
concluyente de superioridad entre algoritmos.
